
\clearpage
% ============================================================================
\section*{v4.10.0-r74 Source-Only: The Value of the Lepton Angle Is Not Fixed by Asymptote, Frame Map, Expansion, Locking or a Sector-Exchange Breaker; the Offset Is Carried by \texorpdfstring{\(\varphi\)}{phi}, the Coefficient \texorpdfstring{\(q\)}{q} Only by Its Sign; the Strand Rests (H228--H236, k442--k451)}
\addcontentsline{toc}{section}{v4.10.0-r74 Source-Only: The Value of the Lepton Angle Is Not Fixed by Asymptote, Frame Map, Expansion, Locking or a Sector-Exchange Breaker; the Offset Is Carried by \texorpdfstring{\(\varphi\)}{phi}, the Coefficient \texorpdfstring{\(q\)}{q} Only by Its Sign; the Strand Rests (H228--H236, k442--k451)}
\label{sec:v4100-r74-phi-strand-rests}
% ============================================================================

\paragraph{Purpose.}
Layer r73 left the common lepton offset as a dimensionless rotation angle \(\varphi\) with the measured value \(\varphi=0.0395773\), the branch
of the data chosen by the sign of the cubic spectral invariant, and that sign set by the author.  On 28 November 2026 nine further tests
followed (H228--H236, scripts k442--k448, k450, k451), each fixed in an advance note before the computation, together with one exploration
without advance note (k449).  This layer records them, on the author's decision of 28 November 2026 to prepare the release.  It adds no new
computation.  The outcome in one sentence: no route tried fixes the value of \(\varphi\) (asymptote of a loop or of the flow, loosening of \(J\)
at order six, transfer of the symmetry point to the matter frame, expansion as a clock, backward search for a breaker of the sector exchange,
locking to eight positions, a common fine step, a Casimir twist between an inner and an outer field); at every equilibrium \(\varphi\) is
exactly the mean of the two sector offsets; and the coefficient \(q\) of the sixth-order term of H167 carries only the sign of the branch choice
(\(q\to0^-\)), while \(\varphi\) carries the offset, which leaves every earlier verdict and the recorded neutrino predictions of r73.10
unchanged.  On 28 November 2026 (15:41) the author decided to leave \(\varphi\) open and to let the \(\varphi\) strand rest; H235 is the one
additional attempt released with the exploration k449, H236 a clarification of the bookkeeping.  Names follow r69: rest frame R (formerly
reference frame), matter frame M, dual frame T (formerly tachyon frame), light point P (formerly photon eigenframe); Im sector (charged
leptons, \(c=1\)) and Re sector (neutral leptons, \(c=0\)), under H99.  The tests H173--H212 remain documented in the research notes and are not
recorded as a layer.  Earlier layers are not rewritten.

\noindent\textbf{Source-only status.}
This v4.10.0-r74 layer is source-only.  It does not build a PDF, does not assign a new DOI,
and does not authorize a public release.  The v5.5.0 Version DOI
\href{https://doi.org/10.5281/zenodo.24563811}{10.5281/zenodo.24563811} and the Concept DOI
\href{https://doi.org/10.5281/zenodo.19210728}{10.5281/zenodo.19210728} are unchanged by it.

\subsection*{v4.10.0-r74 setting}
\label{subsec:v4100-r74-setting}

As in r73: Brannen form \(\sqrt{m_k}=\mu_s\bigl(1+\sqrt2\cos(\delta_s+2\pi k/3)\bigr)\) with \(K=\tfrac23\) under \(J\); tree model of H225,
\(V_{\rm Im}=-\cos2(\delta-\pi/4-\varphi)+q\cos3\delta\), \(V_{\rm Re}=\cos2(\delta-\varphi)+q\cos3\delta\), common vacuum energy
\(E(\varphi)=\sum_s\min_\delta V_s\), exact sector exchange \(V_{\rm Re}(-\delta;\pi/4-\varphi)=V_{\rm Im}(\delta;\varphi)\) (r73.8), symmetric
minimum \(\varphi^*=\pi/8\).  Measured: \(x=0.0395773\) (Im sector, pole masses), \(x_\nu=0.04336\pm0.00213\) (Re sector, from \(R\)).
H228--H235 still carried the fit value \(q=-0.0852\) of H167 next to \(\varphi\) where the tree model was evaluated at finite \(q\); r74.10 settles
this.  Tools: Python for every main computation, Wolfram Language independently for every test and for k449; all common numbers agree.
Qiskit was not used (no circuit).  No HDL/MMA comparison: no finite search or selection procedure of \texttt{RUN\_TAPE} is involved.

\subsection*{v4.10.0-r74 results}
\label{subsec:v4100-r74-results}

\paragraph{Theoremlet v4.10.0-r74.1 (\(\varphi\) is neither the limit of a simple loop nor a natural stage of the flow; H228, k442; refuted).}
(A) Among \(111\,110\) words of length \(1\) to \(5\) in ten angle operations of the frame (\(46\,993\) attracting fixed points in \((0,\pi/4)\))
none lies within \(x\pm2\sigma_\varphi\) (\(\pm1.56\cdot10^{-6}\)); the nearest is \(27\sigma\) off, and the neighbourhood of \(x\) is thinner than that of
a typical control value (\(77\%\) of the controls have at least as many fixed points nearby).  Post hoc: at this precision a loop search of length
five is near its limit of resolution (random expectation \(0.39\) hits in the window).  (B, not blind) In leading order the gradient flow of
\(E\) from \(0\) to \(\pi/8\) has the closed form \(\varphi(t)=\pi/8-\tfrac23\arctan\bigl(\tan(3\pi/16)\,e^{-t/\tau}\bigr)\), \(\tau=1/E''(\pi/8)\)
(Wolfram DSolve and NDSolve); the data value is reached at \(t_x/\tau=0.132105\).  None of ten clocks fixed in advance (\(\tfrac18\),
\(\beta_{\rm sym}^2\), \(1/\gamma_L^2\), \(\tfrac13\), \(1\), \(\tfrac12\), \(\ln2\), \(\varepsilon\), \(\beta_{\rm sym}\), \(1/\gamma_L\)) hits; the nearest,
\(\tfrac18\), lies \(2652\sigma\) off.  Reading \(\varphi\) as an asymptote only moves the measured number \(\varphi\) into a measured number
\(t_x/\tau\).

\paragraph{Theoremlet v4.10.0-r74.2 (loosening \(J\) at order six chooses the branch but not the angle; H229, k443; refuted for \(\varphi\)).}
On the author's choice the full square \(\lambda N^2\) (\(\lambda\ge0\)) of the cubic norm was admitted on trial; at \(K=\tfrac23\) its
\(\delta\)-part is \(2r^6g(\delta)\), \(g=-2\sqrt2\cos3\delta+\cos6\delta\) (Wolfram symbolic).  (a, not blind) For all strengths
\(w=10^{-4}\ldots3\) and \(\varphi=0.01\ldots0.06\) both sectors choose the data branch: stability supplies the sign that r73.6 had to set.
(b) The global minima of \(E(\varphi;w)\) are only the pair \(\pi/12,\pi/6\) (\(w\le0.03\)) or \(\pi/8\) (\(w\ge0.1\)); no \(w\) gives \(x\).
(c) Without rotation the Im data branch is a stationary point of \(g\) (\(g'(5\pi/4)=0\)), so \(N^2\) shifts only the Re branch; the nearly equal
offsets of the data do not come from one strength.  The loosening exchanges a set sign for a changed rule.  It is not adopted: the sign of
the branch-choice term remains the setting of r73.6 (author).  A first version of part (c) used the wrong representative of the Im branch; it
was corrected before part (c) was evaluated, and the discarded value is noted in the output file.

\paragraph{Theoremlet v4.10.0-r74.3 (no frame map carries the symmetry point \(\pi/8\) to the measured value; H230, k444; refuted).}
The symmetry point \(\pi/8\) of \(E\) lies in the rest frame R; \(x\) is measured in the matter frame M.  The missing transfer was made.
(A) The mask on the spectrum with forms K, linear and S at \(\pm\varepsilon_0\), \(\pm2\varepsilon_0\): form S (similarity) leaves the phase
unchanged in every frame (\(\pi/24\) in both sectors); forms K and linear destroy \(K=\tfrac23\) (\(K'=0.744\) to \(23.7\)) and nowhere give the
data phase.  (B) Aberration, contraction and mask with \(\beta_{\rm sym}\) or \(\tanh2\varepsilon_0\), readings \(\varphi,2\varphi,3\varphi\):
the \(36\) images lie between \(0.116\) and \(1.200\); the point \(0\) is fixed by all of them.  Limit cases: the light point P sends angles to
\(0\) or \(\pi\), the dual frame T (\(z\mapsto iz\)) shifts by \(\pi/2\).  Within this frame there is no frame-dependent phase at fixed Koide
value.  The value of \(\varphi\) is not a problem of translation between R and M; the step is closed.

\paragraph{Theoremlet v4.10.0-r74.4 (expansion as a clock: Hubble friction does not freeze \(\varphi\) at the measured value; H231, k445; refuted).}
Set up (in the advance note) as a field rolling in the potential \(E\) from \(0\) in a flat \(\Lambda\)CDM background with Planck values,
\(\ddot\varphi+3H\dot\varphi+E'(\varphi)=0\), without back reaction.  For \(m\ll H_0\) the field stays at its starting value; for \(m\gg H_0\) it rolls
early and oscillates about \(\pi/8\).  Exactly one mass, \(\mu^*=m/H_0=0.9721\), gives \(\varphi=x\) today, and there the field is rolling
(\(d\varphi/d\ln a=0.075\)).  Since \(d\ln m_e/d\varphi=51.5\), this gives \(|d\ln m_e/dt|=2.7\cdot10^{-10}\) per year, about \(2.7\cdot10^6\) times
the atomic-clock bound of order \(10^{-16}\) per year (order of magnitude from the literature, not recomputed).  Natural values of \(\mu\)
(\(\sqrt{\Omega_\Lambda}\), \(1\), \(\beta_{\rm sym}\), \(\tfrac32\), \(\gamma_L\), \(1/\beta_{\rm sym}\)) do not hit and carry the same drift.  A frozen
\(\varphi=x\) is an initial condition.

\paragraph{Theoremlet v4.10.0-r74.5 (backward search for the breaker: \(\varphi\) is the mean of the sector offsets; a symmetric term makes \(x\) only a local minimum; H232, k446; refuted).}
(A1, not blind, Wolfram symbolic) \(dE/d\varphi=-4\cos(o_{\rm Im}-o_{\rm Re})\sin(o_{\rm Im}+o_{\rm Re}-2\varphi)\); hence for every additional term that
depends on \(\delta\) only, at every stationary \(\varphi\),
\[
o_{\rm Im}+o_{\rm Re}=2\varphi ,
\]
exact.  The data then require \(\varphi^*=(x+x_\nu)/2=0.04147\pm0.00106\), and the distance \(x_\nu-x\) (\(1.78\sigma\), r73.1) is twice the shift of
the offsets against \(\varphi\).  (B1) A symmetric spectral term \(q\cos3\delta+r\cos6\delta\) at \(q=-0.0852\) makes \(x\) stationary at
\(r^*=0.0450231\) and a minimum in \(\delta\), but not the global one (\(E(7\pi/24)\) lies lower), and gives \(x_{\nu,\rm pred}=0.09185\), \(+22.8\sigma\).
(B2) The leading-order ratio \(r/|q|=\tfrac23\) does not survive at finite \(q\) (\(0.528\)).  (C, report) For no \(q\) is \(x\) global; a term of
the size of the fine-structure constant in the Im sector alone does not reach \(x\).  As in H226, \(x\) becomes at most a local, stable minimum.

\paragraph{Theoremlet v4.10.0-r74.6 (the fine angles do not lock to eight positions; H233, k447; refuted).}
On the author's picture ``eight positions, each bit carrying twice as much'': grids \((\pi/4)2^{-k}\) and \((\pi/12)2^{-k}\), \(k=0,\ldots,8\).
The tree values \(\pi/12\), \(\pi/6\), \(\pi/4\) and the eighth turns lie on the grids (known, H167, H171).  \(\varphi\) does not: after eight halvings
it is still \(3.06\cdot10^{-4}\) rad, i.e.\ \(392\sigma\), from the nearest point, and it is no remainder of a natural angle after locking.  Eight
grid hits of the mixing angles (\(\theta_{23}\approx\pi/4\), \(\theta_{12}\approx3\pi/16\), \(\theta_{13}\approx3\pi/64\)) match the random
expectation \(5.7\) (\(P(\ge8)=0.22\)).  Locking fixes the tree values, not the fine angles; \(\varphi\) is a continuous remainder.

\paragraph{Theoremlet v4.10.0-r74.7 (no common fine step; H234, k448; refuted).}
Do the fine values \(\varphi\), \(\beta_{\rm sym}-\tfrac12\), \(D/B-1\), \(x_\nu-x\) follow from one small number, \(F=c\,\epsilon^n\) (15 simple factors,
\(n=1,2,3\))?  With \(\epsilon=\varphi\) the sharp value \(\beta_{\rm sym}-\tfrac12\) misses by \(36\sigma\) (best \(\sqrt2\,\varphi\)); with free \(\epsilon\)
the hits occur as often as chance predicts for this density of formulas (\(p=0.74\)).  Post hoc and not counted: \((\beta_{\rm sym}-\tfrac12)/\varphi
=1.2483\), near \(\tfrac54\).  The distinction ``discrete values derived, continuous values open'' describes the state; it saves no input.

\paragraph{Remark v4.10.0-r74.8 (inner and outer field with a Casimir gap; exploration k449, not blind; no number, no prediction).}
On the author's picture of a separated inner region of the original field, an outer field resisting collapse and a vacuum layer between them.
Computed exactly with two tools: a real field with \(V=\lambda(\phi^2-1)^2\) has the two vacua \(\pm1=S^0\), the unit sphere of HS(0), separated by
a kink (BPS remainder \(0\)); two regions give projectors \(e_\pm\) with \(e_+e_-=0\), \(e_++e_-=1\), and \(j=e_+-e_-\) with \(j^2=+1\), i.e.\ exactly
\(\mathbb R\oplus\mathbb Rj\) = HS(\(-1\)) (consistent with k405); in a gap with modes \((n+a)\pi/d\), equal boundary conditions (DD) attract and mixed
ones (DN) repel, ratio \(-7/8\); a relative phase \(\alpha\) across the gap gives \(E(\alpha)\propto-\sum_n\cos(n\alpha)/n^4\) with \(S(\pi)/S(0)=-7/8\)
(exact via \(\operatorname{Re}\operatorname{Li}_4(e^{i\alpha})\)), critical points only at \(0,\pi,2\pi\).  The ratio \(7/8=1-2^{-3}\) is the known
counting number \(\eta(4)/\zeta(4)\).  \emph{Reading, not derived:} two regions give HS(\(-1\)), merged regions HS(0) \(=\mathbb R\), a continuous phase
HS(1) \(=\mathbb C\); the last is admitted by hand, not generated.  Post hoc: only the repulsive case keeps the separation stable, and it is
the case \(j^2=+1\).  The exploration gives no number; the continuous phase is set; \(\varphi\) stays open.

\paragraph{Theoremlet v4.10.0-r74.9 (the Casimir twist breaks the sector exchange but does not place \(\varphi\) at the measured value; H235, k450; refuted).}
Twist energy from k449 with counting weights \(2\) (Im sector, complex) and \(1\) (Re sector, real), fermionic sign, free strength \(\kappa\), in
the tree model with \(q=-0.0852\).  The weights \(2:1\) break the exchange.  (A) A term that depends on \(\varphi\) only cannot change the neutrino
prediction (the offsets are tree functions): requiring \(o_{\rm Im}=x\) gives \(x_{\nu,\rm pred}=0.06295\), \(+9.20\sigma\).  (B1) \(\kappa^*=0.1056\) makes
\(\varphi_x\) a local minimum; (B2) the global minimum lies near \(2\pi/3\), at the mirror partner with the wrong sign of the Im offset.  Post hoc
(k450c, not counted): with small \(q\) the neutrino prediction moves to \(-0.5\sigma\ldots-1.8\sigma\), \(\varphi_x\) is never global, and the model
counts the offset twice when the fit \(q\) of H167 and \(\varphi\) stand together; this led to r74.10.

\paragraph{Theoremlet v4.10.0-r74.10 (double counting settled: \(\varphi\) carries the offset, \(q\) only its sign; no earlier verdict changes; H236, k451; holds).}
(A) Reading Q (\(q=-0.0852\) carries, not blind): \(x_{\nu,\rm pred}=0.06295\), \(+9.20\sigma\).  Reading \(\Phi\) (\(\varphi\) carries, \(q\to0^-\), not
blind): \(x_{\nu,\rm pred}=x\), \(-1.78\sigma\).  Both free (blind): the \(2\sigma\) range of \(q\) is
\[
q\in[-0.029,\,0],
\]
the end \(-0.029\) corresponding to \(x_\nu+2\sigma\) and the end \(0\) arising because \(x_\nu-2\sigma<x\), which would need \(q>0\), the
branch rejected by r73.5.  So \(q=0\) is admitted and \(-0.0852\) lies three
times further out: \(\varphi\) carries the offset, \(q\) is only the sign of the branch choice.  (B, blind) Part B of H226 repeated in reading
\(\Phi\) (\(q=-0.01\), \(-0.001\)): T1 still gives \(x_{\nu,\rm pred}\approx0.48\) (\(+207\sigma\)), T2 and T3 have no solution; the verdict of H226 stands.
(C) Inventory of H221--H235: no verdict changes; numbers that change are the neutrino distances in H232 and H235 (\(-1.3\sigma\), \(-1.8\sigma\)
instead of \(+22.8\sigma\), \(+9.2\sigma\)) and the clock of H228 (\(0.1321\) in reading \(\Phi\) instead of \(0.1231\)); the remark ``\(\tfrac18\) only
\(1.5\%\) off'' of H228 holds only at \(q=-0.0852\) (\(5.7\%\) in reading \(\Phi\)).  The lepton masses and the recorded predictions of r73.10 (H227:
\(\sum m_\nu=59.30\pm0.24\) meV, \(m_\beta\), \(m_{\beta\beta}\), \(R\)) are not affected.  From now on the tree model is used with \(q\to0^-\) (sign set,
value zero) and the offset in \(\varphi\); a finite \(q\) would have to be justified as a scale part.

\paragraph{Author decision v4.10.0-r74.D (28 November 2026, 15:41).}
\(\varphi=0.0395773\) is left open as a measured input, and the \(\varphi\) strand rests.  H235 was the one additional attempt released with
the exploration k449; H236 settles the bookkeeping and makes no claim on the value of \(\varphi\).  OP3 remains open.

\subsection*{v4.10.0-r74 balance of settings}
\label{subsec:v4100-r74-balance}

\begin{center}
\small
\begin{tabular}{@{}>{\raggedright\arraybackslash}p{0.34\textwidth}>{\raggedright\arraybackslash}p{0.58\textwidth}@{}}
\toprule
item & status\\
\midrule
\(\varphi\) at every equilibrium & derived: exact mean of the two sector offsets for any \(\delta\)-dependent addition (r74.5)\\
a breaker depending on \(\varphi\) only & derived: cannot change the neutrino prediction (r74.9)\\
minima of \(E\) under the sector exchange & derived: only at distinguished points \(\pi/12\), \(\pi/8\), \(\pi/6\) for the \(N^2\) ratio (r74.2)\\
transfer of \(\pi/8\) from R to M & done; no frame map gives \(x\) (r74.3)\\
asymptote, expansion clock, eight-position grid, fine step, Casimir twist & do not fix \(\varphi\) (r74.1, r74.4, r74.6, r74.7, r74.9)\\
sign of the branch-choice term & set by the author (r73.6); loosening of \(J\) at order six would replace it by a rule change, not adopted (r74.2)\\
coefficient \(q\) of H167 & only a sign, \(q\to0^-\); data allow \(|q|\le0.029\) (\(2\sigma\)) and require none (r74.10)\\
value \(\varphi=0.0395773\) & measured; left open by the author, strand rests (r74.D)\\
\(K=\tfrac23\), Brannen form, H99, scales \(\mu_e,\mu_\nu\), circulant sector, tree level & set (inputs), as before\\
HS(\(-1\)), HS(0), HS(1) from inner/outer field & reading of an exploration; the continuous phase for HS(1) set (r74.8)\\
\bottomrule
\end{tabular}
\end{center}
Since r73.4 the open choice sits in the value of \(\varphi\), and this layer did not move it: eight further routes left it where it was.  Four
results point the same way, without deciding it: in H226, H232 and H235 the measured value becomes at most a local, stable minimum, and in H231
a frozen value is the starting value.  The target \(\varphi\) is itself a measured input.  One question arose only in a chosen decomposition:
the tree model carried the offset twice (fit \(q\) and rotation \(\varphi\)); r74.10 removes this by fixing that \(\varphi\) carries it.  The fork
``equilibrium value against initial or boundary condition (metastable vacuum)'' was put to the author, who left \(\varphi\) open (r74.D).

\subsection*{v4.10.0-r74 frames}
\label{subsec:v4100-r74-frames}

\begin{center}
\small
\begin{tabular}{@{}>{\raggedright\arraybackslash}p{0.34\textwidth}>{\raggedright\arraybackslash}p{0.22\textwidth}>{\raggedright\arraybackslash}p{0.36\textwidth}@{}}
\toprule
item & frame & comparison\\
\midrule
r74.1, r74.2, r74.5, r74.6, r74.7, r74.9, r74.10 (phases, potential, grids) & M & pole masses; \(K\) and \(R\) frame independent in form S\\
r74.3 (symmetry point) & R \(\to\) M; P and T as limit cases & images of \(\pi/8\) and \(0\) compared with \(x\) in M\\
r74.4 (expansion clock) & M (cosmological background) & drift of \(m_e\) against atomic clocks\\
r74.8 (inner/outer field) & no frame & no comparison with data\\
\bottomrule
\end{tabular}
\end{center}
The sector exchange (\(\delta\to-\delta\), \(\varphi\to\pi/4-\varphi\), Im \(\leftrightarrow\) Re), the frame exchange \(J\), the dual map
\(z\mapsto iz\) and time reversal \(z\mapsto\bar z\) are kept apart (r51, r81).

\subsection*{v4.10.0-r74 register entries}
\label{subsec:v4100-r74-register}

\begin{center}
\small
\begin{tabular}{@{}>{\raggedright\arraybackslash}p{0.40\textwidth}>{\raggedright\arraybackslash}p{0.52\textwidth}@{}}
\toprule
statement & status after v4.10.0-r74\\
\midrule
H228, H230, H231, H232, H233, H234, H235 & refuted\\
H229 & refuted for \(\varphi\); branch chosen by stability under the loosened \(J\); loosening not adopted\\
H236 & holds: \(\varphi\) carries the offset, \(q\to0^-\), \(2\sigma\) range \(q\in[-0.029,0]\)\\
k449 & exploration (not blind); reading, no number\\
verdicts of H221--H235 & unchanged (r74.10)\\
\(\sum m_\nu=59.30\pm0.24\) meV, \(m_\beta=8.87\) meV, \(m_{\beta\beta}\in[1.29,4.00]\) meV, \(R=0.030839\) (H227) & predictions, open, unaffected\\
\(\varphi=0.0395773\) & measured input; left open (author, 28 November 2026, 15:41)\\
sign of the branch-choice term & setting (author, 28 November 2026), as in r73\\
symmetry point \(\pi/8\) in M & transfer done (H230); no frame effect\\
OP3, OP4 & open\\
\bottomrule
\end{tabular}
\end{center}

\subsection*{v4.10.0-r74 notes}
\label{subsec:v4100-r74-notes}

{\raggedright
\texttt{notes/\allowbreak research/\allowbreak hilbert\_\allowbreak shift\_\allowbreak felder\_\allowbreak 2026\_\allowbreak 11\_\allowbreak 24/}: folders
\texttt{asymptote\_\allowbreak phi\_\allowbreak h228/}, \texttt{j\_\allowbreak ordnung6\_\allowbreak h229/}, \texttt{symmetriepunkt\_\allowbreak frame\_\allowbreak h230/},
\texttt{expansion\_\allowbreak uhr\_\allowbreak h231/}, \texttt{rueckwaerts\_\allowbreak brecher\_\allowbreak h232/}, \texttt{einrastung\_\allowbreak f18\_\allowbreak h233/},
\texttt{feinstufe\_\allowbreak h234/}, \texttt{innen\_\allowbreak aussen\_\allowbreak k449/}, \texttt{casimir\_\allowbreak verdrehung\_\allowbreak h235/},
\texttt{doppelzaehlung\_\allowbreak h236/}: advance notes (none for k449), result notes, scripts \texttt{k442} to \texttt{k451} (each with a Wolfram
cross-check \texttt{b}; post hoc addendum \texttt{k450c}; part B of H236 in \texttt{k451\_\allowbreak teilB\_\allowbreak h226.py}) and their actual outputs.
Layer notes: \texttt{notes/\allowbreak research/\allowbreak v4\_\allowbreak 10\_\allowbreak 0\_\allowbreak r74\_\allowbreak phi\_\allowbreak strang\_\allowbreak ruht/}.
Tools: Python and Wolfram Language for every test; no Qiskit (no circuit); no HDL/MMA comparison (no finite search or selection procedure
of \texttt{RUN\_TAPE} involved).
\par}

\subsection*{v4.10.0-r74 ledger and guard update}
\label{subsec:v4100-r74-ledger-guard}

\[
R_{\rm chk}^{v4100\text{-}r74}=(1,1,1)
\]
for (status recorded; computations checked in the cited scripts with independent Wolfram cross-checks; comparison with data in r74.4, r74.5,
r74.9, r74.10; \(\varphi\), the branch sign and the scales not derived).

The layer spends no guard:
\[
R_{\rm glob}^{(v4100\text{-}r74,+)}=(0,1,1,1,1),
\qquad
\Gamma_{\rm guard}^{v4100\text{-}r74}=(1,1,1,1,1),
\qquad
\beta_4=1/256\ \text{unspent}.
\]
The layer does not imply exact finite \(C5\) zero closure, exact all-readout terminality, \(C6/T2\), a global photon-clock theorem,
or global all-readout closure.  OP1--OP5, C5 and C6/T2 remain open.
